\section{Related Work}

Recent clinical-AI research attacks different layers of auditability. MedLog proposes a standardized event-level logging protocol intended to support monitoring and post-deployment analysis across single-shot and agentic clinical AI workflows \citep{Noori.2025}. DoAtlas instead compiles heterogeneous research evidence into executable causal estimands with explicit contrasts, uncertainty, conflict handling, and validation signals \citep{Li.2026}. HEG-TKG focuses on evidence-traceable temporal knowledge graphs, programmatically verified citations, and temporal grounding for rare-disease reasoning \citep{Ahmed.2026}.

Evidence-Graded Decision Authorization (EGDA) makes a particularly important distinction between information supply and inferential governance: retrieved evidence does not by itself determine what a model should be permitted to assert \citep{Lin.2026}. MedScale's authority separation is evaluated against this conceptual prior rather than claiming that governed assertions are unique to MedScale. The difference under study is scope: MedScale attempts to carry authority distinctions beyond model output into source/derived artifacts, human review, compute/integration boundaries, recovery, and external effects.

Operational-safety work likewise argues that predictive accuracy alone is insufficient. SA-ROC converts reliability requirements into operational zones for automation and mandatory human review \citep{Kim.2026}; clinical-AI monitoring reviews document the broader post-deployment monitoring problem \citep{Andersen.2024}; and governance-oriented work emphasizes bounded autonomy, provenance, logging, escalation, and abstention in LLM-enabled clinical workflows \citep{Cao.2026}. FHIRTrustBench further frames interoperability quality as a contributor to clinical-AI trustworthiness rather than a purely infrastructural issue \citep{Bukhari.2026}.

The novelty claim in this paper is therefore deliberately narrow and remains falsifiable. We do not claim to originate provenance, logging, abstention, evidence grading, or auditable inference. We test whether their systems-level complements---authority separation, provenance-bound artifacts, explicit failure semantics, fail-closed effects, and recovery qualification---can be implemented together as inspectable invariants in one local-first research architecture. The source-backed comparison matrix is maintained separately and uses ``not reported'' rather than ``no'' when a paper does not establish a property.
